\subsection{Multilinear Algebra}

Here are two statements about multilinear maps, which are generalizations of the corresponding statements for linear maps.

\begin{lemma}
    \label{lem:continuity_multilinear_maps}
    Let $V_1, \ldots, V_k$ and $W$ be normed vector spaces and 
    \[
    A \colon V_1 \times \cdots \times V_k \to Z
    \]
    be a multilinear map, then $A$ is continous (with respect to the product norm on $V_1 \times \cdots \times V_k$) if and only if it is bounded, i.e. there exists $C>0$ with
    \[
    \lVert A(v_1, \ldots, v_k) \rVert \leq C \prod_{i=1}^{k} \lVert v_i\rVert
    \]
    for all $v_i \in V_i$.
\end{lemma}
\begin{proof}
    The proof is analogous to the linear case. First assume that $A$ is continous, then in particular there exists $\delta >0$, such that 
    \[
    \sum_{i=1}^{k} \lVert v_i \rVert \leq k\delta \quad \Rightarrow \quad \lVert A(v_1, \ldots, v_k) \rVert \leq 1.
    \]
    So for arbitrary $v_i \in V_i$, we have
    \begin{align*}
        \lVert A(v_1, \ldots, v_k) \rVert = \frac{1}{\delta^k} \prod_{i=1}^{k} \lVert v_i\rVert \; \lVert A(\frac{\delta v_1}{\lVert v_1 \rVert}, \ldots, \frac{\delta v_k}{\lVert v_k \rVert}) \rVert \leq \frac{1}{\delta^k} \prod_{i=1}^{k} \lVert v_i\rVert
    \end{align*}
    Now assume, that $A$ is bounded. Given any $v_i,w_i \in V_i$ we compute:
    \begin{align*}
        \lVert A(v_1, \ldots, v_k) - A(w_1, \ldots, &w_k) \rVert  = \lVert \sum_{i=1}^{k} A(w_1,\ldots,w_{i-1}, v_i - w_i, v_{i+1},\ldots,v_k) \rVert \\ 
        &\leq C \sum_{i=1}^{k} \lVert w_1\rVert \cdots \lVert w_{i-1}\rVert \, \lVert v_i - w_i \rVert \, \lVert v_{i+1}\rVert \cdots \lVert v_k \rVert
    \end{align*}
    Since the left hand side vanishes for $v_i \to w_i$ the claim follows.
\end{proof}

\begin{lemma}
    Multilinear maps on finite dimensional vector spaces are continous (with respect to any norm).
\end{lemma}
\begin{proof}
    For ease of notation we only proof the Lemma for bilinear maps, but the proof is completely analogous for the general case. So let $A \colon V \times W \to Z$ be a bilinear map for finite dimensional vetcor spaces $V, W$ and let $e_1, \ldots e_k$ and $d_1,\ldots,d_m$ be a basis for $V$ and $W$, respectively. Given $v \in V$ and $w \in W$ write 
    \[
    v = \sum_{i = 1}^{k} v^i e_i, \; w = \sum_{j = 1}^{m} w^j d_j.
    \]
    Define $C := \mathrm{max}_{i,j} \lVert A(e_i,d_j) \rVert$, then
    \begin{align*}
        \lVert A(v,w) \rVert \leq \sum_{i,j} \lvert v^i \rvert \lvert w^j \rvert \, \lVert A(e_i,d_j) \rVert \leq C \sum_{i,j} \lvert v^i \rvert \lvert w^j \rvert = C \left(\sum_i \lvert v^i \rvert\right) \left(\sum_j \lvert w^j \rvert\right).
    \end{align*}
    Note that 
    \[
    \lVert v \rVert := \sum_i \lvert v^i \rvert
    \]
    defines a norm, so $A$ is bounded with respect to these norms on $V,W$, hence continous by Lemma \ref{lem:continuity_multilinear_maps}. Since all norms on finite dimensional vector spaces are equivalent, the statement follows.
\end{proof}

The following lemma extends the universal property of the Tensor Product to multilinear maps.

\begin{lemma}
    For $i = 1, \ldots n$ let $V_1^i, \ldots V_{k_i}^i$ be vector spaces and
    \[
    \pi_i \colon V_1^i \times \ldots \times V_{k_i}^i \to V_1^i \otimes \ldots \otimes V_{k_i}^i
    \]
    the canonical product maps. Then for every multilinear map $L \colon \prod_{i=1}^n V_1^i \times \ldots \times V_{k_i}^i \to Z$ to a vector space $Z$, there exists a unique multilinear map $A \colon \prod_{i=1}^n V_1^i \otimes \ldots \otimes V_{k_i}^i \to Z$
    such that the following diagram commutes
    \[\begin{tikzcd}[column sep = 4em]
	{\prod_{i=1}^n V_1^i \times \ldots \times V_{k_i}^i} & Z \\
	{\prod_{i=1}^n V_1^i \otimes \ldots \otimes V_{k_i}^i}
	\arrow["L", from=1-1, to=1-2]
	\arrow["{\prod_i\pi_i}"', from=1-1, to=2-1]
	\arrow["A"', dashed, end anchor = {[xshift = 5, yshift = -3]}, from=2-1, to=1-2]
\end{tikzcd}\]
\end{lemma}

\begin{proof}
    Consider the following commutative diagram
    \[\begin{tikzcd}
	{\prod_{i=1}^n V_1^i \times \ldots \times V_{k_i}^i} \\
	{\bigotimes_{i=1}^n V_1^i \otimes \ldots \otimes V_{k_i}^i} && Z \\
	{\prod_{i=1}^n V_1^i \otimes \ldots \otimes V_{k_i}^i}
	\arrow["\pi"', from=1-1, to=2-1]
	\arrow["L", shift left=2, shorten <=6pt, from=1-1, to=2-3]
	\arrow["{\prod_i\pi_i}"', shift right=12, bend right = 55, from=1-1, to=3-1]
	\arrow["B", dashed, from=2-1, to=2-3]
	\arrow["p", from=3-1, to=2-1]
	\arrow["A", shift right=2, shorten <=6pt, from=3-1, to=2-3]
\end{tikzcd}\]
where $p$ and $\pi$ are the canonical product maps, the map $B$ is induced by $L$ by the universal property of the tensor product and $A$ is defined as the composition of $B$ and $p$. Clearly $A$ is the unique map satisfying the statement of the lemma.
\end{proof}